% GENERATED FILE. Edit canonical lessons, then run npm run generate:books.
% canonical-source-hash: fnv1a64:c142ea1d2a0a0641
% canonical-lessons: AR-C392-muhadir, AR-C392-ustadh, AR-C392-ashrafa, AR-C392-mushrif, AR-C392-dhakara

\chapter{Lecturers and supervisors}
\label{ch:ar-lecturers-and-supervisors}
\hlchaptermodality{392}

\begin{chapteropening}
\textbf{By the end of this chapter:} \emph{I can talk about a lecturer, a professor and a supervisor, and say that someone supervised or studied.}

Complete the last lesson of chapter 392: I can talk about a lecturer, a professor and a supervisor, and say that someone supervised or studied.
\end{chapteropening}

\section[\texorpdfstring{\ar{محاضر} (muḥāḍir)}{محاضر (muḥāḍir)}]{\ar{محاضر} (muḥāḍir) — a lecturer}

\label{lesson:AR-C392-muhadir}



\begin{quote}
Before the new one: say the Arabic for an institute, then the Arabic for a questionnaire.
\end{quote}

\subsection*{You'll want to know: \ar{محاضر}}
\textbf{\ar{محاضر}} — \emph{muḥāḍir} — \textquotedblleft{}a lecturer\textquotedblright{}.

\begin{quote}
\ar{يشرح} \ar{المحاضر} \ar{الدرس} \ar{بوضوح}.
\end{quote}

\begin{quote}
\emph{yashraḥu al-muḥāḍir ad-dars bi-wuḍūḥ.}
\end{quote}

\begin{quote}
The lecturer explains the lesson clearly.
\end{quote}

\begin{cousinweb}
\textbf{\ar{محاضر}} (\emph{muḥāḍir}) is built on the native Arabic root \emph{ḥ-ḍ-r}, \textquotedblleft{}to be present\textquotedblright{} (Lane). \emph{Muḥāḍir} is the Form III active participle, one who speaks to those present: a lecturer.
\end{cousinweb}

\begin{grammarlens}[title={What you've built}]
One more word for giving reasons and opinions.
\end{grammarlens}

\subsection*{Your turn}
\begin{itemize}
\raggedright
  \item \emph{Say it:} \emph{muḥāḍir}
  \item \emph{Say it:} \emph{muḥāḍir}, once more
  \item \emph{Say it:} \emph{yashraḥu al-muḥāḍir ad-dars bi-wuḍūḥ.} — The lecturer explains the lesson clearly.
  \item \emph{Recall:} say the Arabic for an institute, then the Arabic for a questionnaire, then say \emph{muḥāḍir} again
\end{itemize}

\begin{tcolorbox}[breakable,colback=teal!4,colframe=teal!35!black,title={Before you move on}]
What does \ar{محاضر} mean? (\textquotedblleft{}A lecturer\textquotedblright{}.) Say it once more.
\end{tcolorbox}
\section[\texorpdfstring{\ar{أستاذ} (ustādh)}{أستاذ (ustādh)}]{\ar{أستاذ} (ustādh) — a professor}

\label{lesson:AR-C392-ustadh}



\begin{quote}
Before the new one: say the Arabic for a questionnaire, then the Arabic for a lecturer.
\end{quote}

\subsection*{You'll want to know: \ar{أستاذ}}
\textbf{\ar{أستاذ}} — \emph{ustādh} — \textquotedblleft{}a professor\textquotedblright{}.

\begin{quote}
\ar{أستاذ} \ar{الأدب} \ar{رجل} \ar{لطيف}.
\end{quote}

\begin{quote}
\emph{ustādh al-adab rajul laṭīf.}
\end{quote}

\begin{quote}
The literature professor is a kind man.
\end{quote}

Plural: \ar{أساتذة}, \emph{asātidha}.

\begin{cousinweb}
\textbf{\ar{أستاذ}} (\emph{ustādh}) was borrowed from Persian \emph{ustād}, \textquotedblleft{}a master, a teacher\textquotedblright{} (EtymDB). In Arabic it is both a polite title and a university professor.
\end{cousinweb}

\begin{grammarlens}[title={What you've built}]
One more word for giving reasons and opinions.
\end{grammarlens}

\subsection*{Your turn}
\begin{itemize}
\raggedright
  \item \emph{Say it:} \emph{ustādh}
  \item \emph{Say it:} \emph{ustādh}, once more
  \item \emph{Say it:} \emph{ustādh al-adab rajul laṭīf.} — The literature professor is a kind man.
  \item \emph{Recall:} say the Arabic for a questionnaire, then the Arabic for a lecturer, then say \emph{ustādh} again
\end{itemize}

\begin{tcolorbox}[breakable,colback=teal!4,colframe=teal!35!black,title={Before you move on}]
What does \ar{أستاذ} mean? (\textquotedblleft{}A professor\textquotedblright{}.) Say it once more.
\end{tcolorbox}
\section[\texorpdfstring{\ar{أشرف} (ashrafa)}{أشرف (ashrafa)}]{\ar{أشرف} (ashrafa) — to supervise}

\label{lesson:AR-C392-ashrafa}



\begin{quote}
Before the new one: say the Arabic for a lecturer, then the Arabic for a professor.
\end{quote}

\subsection*{You'll want to know: \ar{أشرف}}
\textbf{\ar{أشرف}} — \emph{ashrafa} — \textquotedblleft{}to supervise\textquotedblright{}.

\begin{quote}
\ar{أشرف} \ar{المعلم} \ar{على} \ar{الرحلة}.
\end{quote}

\begin{quote}
\emph{ashrafa al-muʿallim ʿalā ar-riḥla.}
\end{quote}

\begin{quote}
The teacher supervised the trip.
\end{quote}

Present: \ar{يشرف}, \emph{yushrif}.

It takes \ar{على}: \ar{أشرف} \ar{على} ….

\begin{cousinweb}
\textbf{\ar{أشرف}} (\emph{ashrafa}) is built on the native Arabic root \emph{sh-r-f}, \textquotedblleft{}to be high, noble\textquotedblright{} (Lane). It is Form IV, to look down from a height: to oversee, to supervise.
\end{cousinweb}

\begin{grammarlens}[title={What you've built}]
One more word for giving reasons and opinions.
\end{grammarlens}

\subsection*{Your turn}
\begin{itemize}
\raggedright
  \item \emph{Say it:} \emph{ashrafa}
  \item \emph{Say it:} \emph{ashrafa}, once more
  \item \emph{Say it:} \emph{ashrafa al-muʿallim ʿalā ar-riḥla.} — The teacher supervised the trip.
  \item \emph{Recall:} say the Arabic for a lecturer, then the Arabic for a professor, then say \emph{ashrafa} again
\end{itemize}

\begin{tcolorbox}[breakable,colback=teal!4,colframe=teal!35!black,title={Before you move on}]
What does \ar{أشرف} mean? (\textquotedblleft{}To supervise\textquotedblright{}.) Say it once more.
\end{tcolorbox}
\section[\texorpdfstring{\ar{مشرف} (mushrif)}{مشرف (mushrif)}]{\ar{مشرف} (mushrif) — a supervisor}

\label{lesson:AR-C392-mushrif}



\begin{quote}
Before the new one: say the Arabic for a professor, then the Arabic for to supervise.
\end{quote}

\subsection*{You'll want to know: \ar{مشرف}}
\textbf{\ar{مشرف}} — \emph{mushrif} — \textquotedblleft{}a supervisor\textquotedblright{}.

\begin{quote}
\ar{ساعدني} \ar{المشرف} \ar{في} \ar{بحثي}.
\end{quote}

\begin{quote}
\emph{sāʿadanī al-mushrif fī baḥthī.}
\end{quote}

\begin{quote}
The supervisor helped me with my research.
\end{quote}

\begin{cousinweb}
\textbf{\ar{مشرف}} (\emph{mushrif}) is built on the native Arabic root \emph{sh-r-f}, \textquotedblleft{}to be high, noble\textquotedblright{} (Lane). \emph{Mushrif} is the Form IV active participle: one who oversees, a supervisor.
\end{cousinweb}

\begin{grammarlens}[title={What you've built}]
One more word for giving reasons and opinions.
\end{grammarlens}

\subsection*{Your turn}
\begin{itemize}
\raggedright
  \item \emph{Say it:} \emph{mushrif}
  \item \emph{Say it:} \emph{mushrif}, once more
  \item \emph{Say it:} \emph{sāʿadanī al-mushrif fī baḥthī.} — The supervisor helped me with my research.
  \item \emph{Recall:} say the Arabic for a professor, then the Arabic for to supervise, then say \emph{mushrif} again
\end{itemize}

\begin{tcolorbox}[breakable,colback=teal!4,colframe=teal!35!black,title={Before you move on}]
What does \ar{مشرف} mean? (\textquotedblleft{}A supervisor\textquotedblright{}.) Say it once more.
\end{tcolorbox}
\section[\texorpdfstring{\ar{ذاكر} (dhākara)}{ذاكر (dhākara)}]{\ar{ذاكر} (dhākara) — to study for}

\label{lesson:AR-C392-dhakara}



\begin{quote}
Before the new one: say the Arabic for to supervise, then the Arabic for a supervisor.
\end{quote}

\subsection*{You'll want to know: \ar{ذاكر}}
\textbf{\ar{ذاكر}} — \emph{dhākara} — \textquotedblleft{}to study for\textquotedblright{}.

\begin{quote}
\ar{ذاكر} \ar{الطالب} \ar{دروسه} \ar{حتى} \ar{منتصف} \ar{الليل}.
\end{quote}

\begin{quote}
\emph{dhākara aṭ-ṭālib durūsahu ḥattā muntaṣaf al-layl.}
\end{quote}

\begin{quote}
The student studied his lessons until midnight.
\end{quote}

Present: \ar{يذاكر}, \emph{yudhākir}.

Common in spoken Arabic; formal writing often says \ar{درس}.

\begin{cousinweb}
\textbf{\ar{ذاكر}} (\emph{dhākara}) is built on the native Arabic root \emph{dh-k-r}, usually traced to the sense \textquotedblleft{}to remember, to mention\textquotedblright{}. It is Form III, to go over a thing with oneself again: to study for a test.
\end{cousinweb}

\begin{grammarlens}[title={What you've built}]
One more word for giving reasons and opinions.
\end{grammarlens}

\subsection*{Your turn}
\begin{itemize}
\raggedright
  \item \emph{Say it:} \emph{dhākara}
  \item \emph{Say it:} \emph{dhākara}, once more
  \item \emph{Say it:} \emph{dhākara aṭ-ṭālib durūsahu ḥattā muntaṣaf al-layl.} — The student studied his lessons until midnight.
  \item \emph{Recall:} say the Arabic for to supervise, then the Arabic for a supervisor, then say \emph{dhākara} again
\end{itemize}

\begin{tcolorbox}[breakable,colback=teal!4,colframe=teal!35!black,title={Before you move on}]
What does \ar{ذاكر} mean? (\textquotedblleft{}To study for\textquotedblright{}.) Say it once more.
\end{tcolorbox}
